\documentclass[11pt,a4paper]{article}

\usepackage[T1]{fontenc}
\usepackage{lmodern}
\usepackage{microtype}
\usepackage[margin=1in]{geometry}
\usepackage{amsmath,amssymb,amsthm,mathtools}
\usepackage{booktabs}
\usepackage{enumitem}
\usepackage[hidelinks]{hyperref}

\theoremstyle{plain}
\newtheorem{theorem}{Theorem}[section]
\newtheorem{lemma}[theorem]{Lemma}
\newtheorem{proposition}[theorem]{Proposition}
\newtheorem{corollary}[theorem]{Corollary}

\theoremstyle{definition}
\newtheorem{definition}[theorem]{Definition}

\theoremstyle{remark}
\newtheorem{remark}[theorem]{Remark}

\newcommand{\R}{\mathbb{R}}
\newcommand{\N}{\mathbb{N}}
\newcommand{\Jcost}{J}

\title{\textbf{Algorithmic Cost and the Halting Problem}\\[0.5em]
\large Why the Universe Cannot Loop:\\
Computation, Divergence, and Economic Censorship\\
from the $J$-Cost Functional}
\author{Jonathan Washburn\\
\small Recognition Science Research Institute, Austin, Texas\\
\small \texttt{jon@recognitionphysics.org}}
\date{March 2026}

\begin{document}
\maketitle

\begin{abstract}
We prove that \emph{every physical computation is bounded by its initial $\Jcost$-cost budget}, that \emph{infinite loops are economically impossible}, and that \emph{the halting problem for ledger-realized computation is resolved by cost exhaustion}.

The argument requires only three ingredients already present in Recognition Science:
(i)~the non-negativity of the defect functional $\Jcost(x) = \tfrac{1}{2}(x + x^{-1}) - 1 \ge 0$,
(ii)~the variational dynamics in which total defect is monotone non-increasing,
and (iii)~the Archimedean property of the reals.

The central result (\textbf{Infinite Computation Impossibility Theorem}, \S\ref{sec:impossible}) states:

\begin{quote}
\itshape
No variational trajectory on the Recognition Science ledger can sustain infinitely many steps each costing at least $\delta > 0$ in defect reduction. After at most $\lfloor D_0/\delta \rfloor$ such steps, where $D_0$ is the initial total defect, the computation must slow below threshold.
\end{quote}

This result unifies two previously separate censorship phenomena.
\texttt{LogicFromCost.lean} proved that contradictions carry positive cost and are therefore censored.
\texttt{AlgorithmicCost.lean} proves that infinite loops carry infinite cumulative cost and are therefore censored.
Both are instances of the same meta-principle: the universe forbids infinite defect accumulation.
A contradiction is a \emph{static} infinite-cost configuration; an infinite loop is a \emph{dynamic} infinite-cost trajectory.
The cost functional censors both.

All theorems are machine-verified in Lean~4 with no \texttt{sorry} and no custom axioms.
The formalization is available at \texttt{RecognitionScience/Foundation/AlgorithmicCost.lean}.
\end{abstract}

\tableofcontents
\newpage

%% ============================================================
\section{Introduction}\label{sec:intro}
%% ============================================================

Recognition Science derives the structure of physical reality from a single cost functional
\begin{equation}\label{eq:jcost}
\Jcost(x) \;=\; \frac{1}{2}\!\left(x + x^{-1}\right) - 1, \qquad x > 0,
\end{equation}
together with the combinatorics of the $3$-cube $Q_3$ and the self-similarity equation $x^2 = x + 1$ that forces $x = \varphi$, the golden ratio.

Previous work established that this framework yields:
\begin{itemize}[nosep]
\item The law of existence: $x$ exists iff $\Jcost(x) = 0$ iff $x = 1$ \cite{RS-LawOfExistence}.
\item Logic from cost: consistency minimizes $\Jcost$; contradictions carry positive cost \cite{RS-LogicFromCost}.
\item Variational dynamics: the ledger evolves by constrained $\Jcost$-minimization, producing deterministic, defect-monotone trajectories \cite{RS-VariationalDynamics}.
\item Temperature, entropy, and the three laws of thermodynamics from the ledger's cost structure \cite{RS-Thermodynamics}.
\end{itemize}

But logic is only the \emph{foundation} of computation, and the theory was silent on a deeper question: \textbf{what are the computational limits imposed by $\Jcost$-cost?}

This paper answers that question.
We prove that every physical computation---modeled as a variational trajectory on the ledger---is bounded by its initial defect budget, that infinite loops are economically impossible, and that the halting problem for physically realized computation is resolved by cost exhaustion.

The argument is elementary.
It uses only the non-negativity and monotonicity of total defect together with the Archimedean property of the reals.
No computability theory, no diagonalization, and no reference to Turing machines is needed.
The universe does not need a halting oracle; it \emph{is} the halting oracle, and its mechanism is $\Jcost$-cost exhaustion.

%% ============================================================
\section{Preliminaries}\label{sec:prelim}
%% ============================================================

\subsection{The Cost Functional}

The canonical cost functional $\Jcost : \R_{>0} \to \R$ is defined by~\eqref{eq:jcost}.
Its key properties, all machine-verified:

\begin{proposition}[Properties of $\Jcost$]\label{prop:jcost}
\begin{enumerate}[nosep,label=(\roman*)]
\item $\Jcost(x) \ge 0$ for all $x > 0$ \quad (AM--GM inequality).
\item $\Jcost(x) = 0 \iff x = 1$ \quad (unique zero).
\item $\Jcost(x) = \Jcost(x^{-1})$ \quad (reciprocal symmetry).
\item $\Jcost$ is strictly convex on $(0,\infty)$.
\item $\Jcost(x) \to +\infty$ as $x \to 0^+$ or $x \to +\infty$.
\end{enumerate}
\end{proposition}

In log-coordinates $t = \log x$, the cost becomes $\Jcost_{\log}(t) = \cosh t - 1$, a convex bowl centered at $t=0$ with $\Jcost_{\log}''(0) = 1$.

\subsection{Ledger Configurations and Total Defect}

A \emph{ledger configuration} of size $N$ is a function $c : \{0,\ldots,N-1\} \to \R_{>0}$ assigning a positive real ratio to each ledger entry.
The \emph{total defect} is
\begin{equation}\label{eq:totaldefect}
D(c) \;=\; \sum_{i=0}^{N-1} \Jcost(c_i) \;\ge\; 0,
\end{equation}
with $D(c) = 0$ iff $c_i = 1$ for all $i$ (the unity configuration).

\subsection{Variational Dynamics}

The ledger evolves by constrained $\Jcost$-minimization.
At each tick $t$, the successor state $c(t{+}1)$ is the configuration minimizing $D$ over the \emph{feasible set}
\[
\mathcal{F}(c) \;=\; \bigl\{c' : D(c') \text{ is defined and } \textstyle\sum_i \log c'_i = \sum_i \log c_i\bigr\},
\]
which preserves the conserved log-charge $\sigma = \sum_i \log c_i$.

A \emph{trajectory} is a sequence $\tau : \N \to \text{Config}(N)$ satisfying
\[
\tau(t{+}1) \;=\; \operatorname*{arg\,min}_{c' \in \mathcal{F}(\tau(t))} D(c')
\qquad \forall\, t \in \N.
\]

\begin{proposition}[Defect Monotonicity]\label{prop:monotone}
Along any variational trajectory, total defect is monotone non-increasing:
\[
D(\tau(t{+}1)) \;\le\; D(\tau(t)) \qquad \forall\, t \in \N.
\]
\end{proposition}

\begin{proof}
The current state $\tau(t)$ is feasible for itself.
The successor minimizes over the feasible set.
Therefore $D(\tau(t{+}1)) \le D(\tau(t))$.
\end{proof}

%% ============================================================
\section{Computation as Trajectory}\label{sec:computation}
%% ============================================================

\begin{definition}[Step Cost]\label{def:stepcost}
The \emph{step cost} of a trajectory $\tau$ at time $t$ is the defect reduction achieved by the variational step:
\[
\delta_t \;=\; D(\tau(t)) - D(\tau(t{+}1)) \;\ge\; 0.
\]
\end{definition}

The step cost is the ``computational work'' performed at tick $t$.
It is non-negative by defect monotonicity.
It is zero iff the state is already at equilibrium---the computation has halted.
It is positive iff genuine work was performed.

\begin{definition}[Cumulative Cost]\label{def:cumcost}
The \emph{cumulative cost} over $T$ steps is
\[
\Delta_T \;=\; D(\tau(0)) - D(\tau(T)) \;=\; \sum_{t=0}^{T-1} \delta_t.
\]
\end{definition}

The telescoping equality follows immediately from the definition.
The cumulative cost is the total computation budget expended from tick~0 to tick~$T$.

\begin{proposition}[Cost Budget]\label{prop:budget}
The cumulative cost is bounded above by the initial defect:
\[
\Delta_T \;=\; D(\tau(0)) - D(\tau(T)) \;\le\; D(\tau(0)) \;=:\; D_0.
\]
\end{proposition}

\begin{proof}
$D(\tau(T)) \ge 0$ by~\eqref{eq:totaldefect}, so $\Delta_T = D_0 - D(\tau(T)) \le D_0$.
\end{proof}

This bound is the key to everything that follows.
The initial defect $D_0$ is the computation's total fuel.
Every step burns some of it.
Once the fuel is exhausted, the computation has halted.

%% ============================================================
\section{The Computation Bound}\label{sec:bound}
%% ============================================================

\begin{theorem}[Step Cost Accumulation]\label{thm:accumulation}
If each of $n$ consecutive steps costs at least $\delta > 0$, the cumulative cost satisfies
\[
\Delta_n \;\ge\; n\delta.
\]
\end{theorem}

\begin{proof}
By induction on $n$.
The base case $n=0$ is trivial ($\Delta_0 = 0 \ge 0$).
For the inductive step, $\Delta_{k+1} = \Delta_k + \delta_{k} \ge k\delta + \delta = (k{+}1)\delta$.
\end{proof}

\begin{theorem}[Computation Bound]\label{thm:bound}
If each step costs at least $\delta > 0$, then the number of such steps satisfies
\begin{equation}\label{eq:bound}
n\delta \;\le\; D_0.
\end{equation}
Equivalently, $n \le D_0/\delta$.
\end{theorem}

\begin{proof}
Combine Theorem~\ref{thm:accumulation} ($\Delta_n \ge n\delta$) with Proposition~\ref{prop:budget} ($\Delta_n \le D_0$).
\end{proof}

\begin{remark}
The bound $n \le D_0/\delta$ is the \emph{computation budget} of the trajectory.
It depends on exactly two quantities: the initial defect $D_0$ (how far the universe is from equilibrium) and the minimum step cost $\delta$ (the granularity of the computation).
No parameter of the computation itself---its algorithm, its data, its purpose---enters the bound.
The budget is a property of the \emph{cost landscape}, not of the computation.
\end{remark}

%% ============================================================
\section{Infinite Loops Are Economically Impossible}\label{sec:impossible}
%% ============================================================

\begin{theorem}[Infinite Computation Impossibility]\label{thm:impossible}
No variational trajectory can sustain infinitely many steps each costing at least $\delta > 0$.
Formally: for any $\delta > 0$,
\[
\neg\;\bigl(\forall\, t \in \N,\; \delta_t \ge \delta\bigr).
\]
\end{theorem}

\begin{proof}
Suppose for contradiction that $\delta_t \ge \delta$ for all $t$.
By the Archimedean property of $\R$, choose $n \in \N$ with $n > D_0/\delta$.
Then by Theorem~\ref{thm:bound}, $n\delta \le D_0$.
But $n > D_0/\delta$ implies $n\delta > D_0$.
Contradiction.
\end{proof}

\begin{corollary}[Eventual Slowdown]\label{cor:slowdown}
For any $\delta > 0$, there exists $T \in \N$ such that $\delta_T < \delta$.
\end{corollary}

\begin{proof}
Contrapositive of Theorem~\ref{thm:impossible}.
If $\delta_t \ge \delta$ for all $t$, then Theorem~\ref{thm:impossible} gives a contradiction.
Therefore $\delta_T < \delta$ for some $T$.
\end{proof}

\begin{corollary}[Asymptotic Halting]\label{cor:asymptotic}
For any $\varepsilon > 0$, the trajectory eventually has a step costing less than $\varepsilon$.
The step costs approach zero in the sense that no positive lower bound can be maintained forever.
\end{corollary}

\begin{remark}
This does not claim that $\delta_t \to 0$ as $t \to \infty$ (which would require a monotone convergence argument).
It claims the weaker but already powerful fact: no positive floor can persist indefinitely.
The universe cannot sustain high-throughput computation forever.
It must eventually relax toward equilibrium.
\end{remark}

%% ============================================================
\section{Equilibrium as Halting}\label{sec:equilibrium}
%% ============================================================

\begin{definition}[Equilibrium / Halting]\label{def:equilibrium}
A configuration $c$ is at \emph{equilibrium} if it minimizes total defect over its feasible set:
\[
\forall\, c' \in \mathcal{F}(c),\quad D(c) \le D(c').
\]
Equivalently, $c$ is a fixed point of the variational dynamics.
\end{definition}

\begin{theorem}[Zero Cost Characterizes Halting]\label{thm:halting-equiv}
A step has zero cost if and only if the state is at equilibrium:
\[
\delta_t = 0 \quad\Longleftrightarrow\quad \tau(t) \text{ is at equilibrium.}
\]
\end{theorem}

\begin{proof}
($\Rightarrow$)
If $\delta_t = 0$, then $D(\tau(t)) = D(\tau(t{+}1))$.
Since $\tau(t{+}1)$ minimizes $D$ over $\mathcal{F}(\tau(t))$ and $\tau(t) \in \mathcal{F}(\tau(t))$, we get $D(\tau(t)) = D(\tau(t{+}1)) \le D(c')$ for all $c' \in \mathcal{F}(\tau(t))$.
So $\tau(t)$ is at equilibrium.

\noindent($\Leftarrow$)
If $\tau(t)$ is at equilibrium, then $D(\tau(t)) \le D(\tau(t{+}1))$ (since $\tau(t{+}1)$ is feasible).
Combined with defect monotonicity $D(\tau(t{+}1)) \le D(\tau(t))$, equality follows.
\end{proof}

%% ============================================================
\section{Loops Imply Fixed Points}\label{sec:loops}
%% ============================================================

\begin{theorem}[No Non-Trivial Loops]\label{thm:noloops}
If a trajectory revisits a previous state---$\tau(t_2) = \tau(t_1)$ with $t_1 < t_2$---then every step in the interval $[t_1, t_2)$ has zero cost, and every state in the interval is at equilibrium.
\end{theorem}

\begin{proof}
The states $\tau(t_1)$ and $\tau(t_2)$ have identical entries, hence identical total defect: $D(\tau(t_1)) = D(\tau(t_2))$.
But defect is monotone non-increasing on $[t_1, t_2]$:
\[
D(\tau(t_1)) \ge D(\tau(t_1{+}1)) \ge \cdots \ge D(\tau(t_2)) = D(\tau(t_1)).
\]
All inequalities are equalities.
Therefore $\delta_t = 0$ for all $t \in [t_1, t_2)$, and by Theorem~\ref{thm:halting-equiv}, each $\tau(t)$ is at equilibrium.
\end{proof}

\begin{remark}
This theorem says that the variational dynamics on the ledger is \emph{loop-free} in a strong sense.
The only ``loop'' is the trivial one: a fixed point that repeats itself forever.
Non-trivial cycles---states that return to a previous configuration without being at equilibrium---are impossible.
An infinite loop would be an infinite repetition of non-trivial states, each costing positive defect.
But that is precisely what Theorem~\ref{thm:impossible} forbids.
\end{remark}

%% ============================================================
\section{Economic Censorship: Unifying Logic and Computation}\label{sec:censorship}
%% ============================================================

The deepest consequence of this work is not any single theorem but the \emph{structural identity} it reveals between two previously separate RS results.

\texttt{LogicFromCost.lean} proved:
\begin{quote}
Contradictions carry positive cost $\Rightarrow$ the universe censors contradictions.
\end{quote}

\texttt{AlgorithmicCost.lean} proves:
\begin{quote}
Infinite loops carry infinite cumulative cost $\Rightarrow$ the universe censors infinite loops.
\end{quote}

Both are instances of the \textbf{Economic Censorship Principle}:

\begin{quote}
\itshape
The universe forbids infinite defect accumulation.
Any configuration or trajectory whose total $\Jcost$-cost would exceed the finite initial budget $D_0$ is physically unrealizable.
\end{quote}

\begin{theorem}[Economic Censorship]\label{thm:censorship}
For any variational trajectory $\tau$ with initial defect $D_0$:
\begin{enumerate}[nosep,label=(\roman*)]
\item $D_0 \ge 0$ \quad (finite, non-negative budget).
\item $D(\tau(t{+}1)) \le D(\tau(t))$ for all $t$ \quad (cost is paid, never refunded).
\item $\Delta_T \le D_0$ for all $T$ \quad (cumulative cost is bounded).
\item For all $\delta > 0$: $\neg(\forall t,\, \delta_t \ge \delta)$ \quad (no infinite loops).
\item For all $\varepsilon > 0$: $\exists T,\, \delta_T < \varepsilon$ \quad (eventual slowdown).
\end{enumerate}
\end{theorem}

The following table summarizes the unification:

\begin{center}
\begin{tabular}{@{}lcc@{}}
\toprule
& \textbf{Contradiction} & \textbf{Infinite Loop} \\
\midrule
Type & Static config & Dynamic trajectory \\
Cost & $\Jcost(P \wedge \neg P) > 0$ & $\sum \delta_t = \infty$ \\
Censorship & Cannot stabilize at $\Jcost = 0$ & Cannot fit in budget $D_0$ \\
Mechanism & Complementarity of ratios & Archimedean property of $\R$ \\
RS theorem & \texttt{contradiction\_positive\_cost} & \texttt{infinite\_computation\_impossible} \\
\bottomrule
\end{tabular}
\end{center}

Both pathologies---logical inconsistency and computational divergence---are the same disease viewed from different angles.
The cost functional cures both.

%% ============================================================
\section{The Halting Theorem for Ledger Computation}\label{sec:halting}
%% ============================================================

\begin{definition}[Computational Process]\label{def:process}
A \emph{computational process} on the ledger is a variational trajectory $\tau$ together with a minimum step cost $\delta > 0$.
The parameter $\delta$ represents the granularity of the computation: the minimum defect reduction per non-trivial step.
\end{definition}

\begin{definition}[Halting Bound]\label{def:haltingbound}
The \emph{halting bound} of a computational process $(\tau, \delta)$ is
\[
T_{\max} \;=\; \frac{D_0}{\delta},
\]
where $D_0 = D(\tau(0))$ is the initial total defect.
\end{definition}

\begin{theorem}[The Halting Theorem]\label{thm:halting}
Every computational process on the ledger reaches effective halting: there exists $T \in \N$ with $T \le T_{\max}$ such that $\delta_T < \delta$.
\end{theorem}

\begin{proof}
Immediate from Corollary~\ref{cor:slowdown} and the explicit bound $n \le D_0/\delta$ (Theorem~\ref{thm:bound}).
\end{proof}

\begin{remark}[The Universe as Halting Oracle]
This theorem does not ``solve'' the abstract Halting Problem, which concerns formal Turing machines with no cost constraint.
What it proves is something different and, in a physical sense, stronger:

\begin{enumerate}[nosep]
\item Abstract Turing machines can loop forever---they have no cost budget.
\item Physical computations (ledger trajectories) cannot---they are bounded by $D_0/\delta$.
\item The universe does not need an external halting oracle.
\item It \emph{is} the halting oracle: its mechanism is $\Jcost$-cost exhaustion.
\end{enumerate}

The halting question is not undecidable for physical computation.
It is decided by the cost landscape itself.
Every physical computation halts because the universe runs out of defect to spend.
\end{remark}

%% ============================================================
\section{Connections to Algorithmic Information Theory}\label{sec:ait}
%% ============================================================

The computation bound $n \le D_0/\delta$ connects naturally to several concepts in algorithmic information theory.

\subsection{Kolmogorov Complexity}

The Kolmogorov complexity of an object is the length of the shortest program that produces it.
In RS terms, the ``program'' is the trajectory from the initial configuration to the target, and the ``length'' is the number of non-trivial steps.
The computation bound says:
\[
K_{\text{RS}}(\text{target}) \;\le\; \frac{D(\text{initial}) - D(\text{target})}{\delta}.
\]
The RS-complexity of a target state is bounded by the defect reduction needed to reach it, divided by the minimum step cost.
This is a cost-theoretic complexity measure that does not depend on a choice of universal Turing machine.

\subsection{Levin's Universal Search}

Levin's universal search allocates computation time to programs in proportion to $2^{-|p|}$, where $|p|$ is program length.
In RS terms, the analogous weight is $e^{-D(c)}$ (the Boltzmann weight from $\Jcost$-cost), which allocates probability to configurations in proportion to their defect.
Lower defect $=$ higher weight $=$ more ``computational probability.''
The Born rule structure proved in \texttt{MeasurementMechanism.lean} is precisely this allocation.

\subsection{The $\varphi$-Connection}

Recognition Science derives all fundamental constants from the golden ratio $\varphi = (1+\sqrt{5})/2$.
The minimum step cost $\delta$ in physical computation is determined by the ledger's discreteness scale, which is itself set by $\varphi$ through the self-similarity equation $\varphi^2 = \varphi + 1$.
Therefore the halting bound becomes
\[
T_{\max} \;\sim\; \frac{D_0}{\delta(\varphi)},
\]
and the computational capacity of the universe is bounded by $\varphi$ and $\Jcost$.

%% ============================================================
\section{Discussion}\label{sec:discussion}
%% ============================================================

\subsection{What This Result Does Not Claim}

We do not claim that the abstract Halting Problem is decidable.
Turing's 1936 theorem is a statement about formal systems without cost constraints.
Our result is about physical computation, which is subject to the $\Jcost$-cost budget.

We do not claim that Turing machines are ``wrong.''
They are correct models of \emph{idealized} computation with unbounded resources.
But the universe is not an idealized computer.
It has finite initial defect, and every transition costs positive defect.
The RS computation bound is a physical constraint, not a logical one.

\subsection{What This Result Does Claim}

\begin{enumerate}[nosep]
\item \textbf{Infinite loops are physically impossible.}
Any process on the ledger that tries to loop forever with positive step cost will exhaust its defect budget in finite time.

\item \textbf{Computation is bounded by cost, not by time or space.}
The fundamental resource for computation is not clock cycles or memory cells but defect---the distance from equilibrium.
A universe with more initial defect can compute more.

\item \textbf{Logic and computation are unified through cost.}
Contradictions (static infinite cost) and infinite loops (dynamic infinite cost) are both censored by the same $\Jcost$-cost bound.
The cost functional provides a single mechanism for both logical consistency and computational boundedness.

\item \textbf{The universe is its own halting oracle.}
No external decider is needed to determine whether a physical computation halts.
The cost landscape itself enforces halting through economic exhaustion.
\end{enumerate}

\subsection{Open Questions}

\begin{enumerate}[nosep]
\item \textbf{Tight bounds.}
The bound $n \le D_0/\delta$ is sharp in principle (it uses only non-negativity and monotonicity).
Can the bound be tightened using the strict convexity of $\Jcost$ or the specific form of the equilibrium (uniform configuration)?

\item \textbf{Computational universality.}
The ledger can ``compute'' in the sense that its trajectories transform configurations.
Is the variational dynamics Turing-complete up to the halting bound?
If so, the ledger can simulate any Turing machine for finitely many steps.

\item \textbf{Complexity classes.}
The bound $T_{\max} = D_0/\delta$ defines a natural complexity class: the set of problems solvable within a given defect budget.
How does this class relate to classical complexity classes (P, NP, PSPACE)?

\item \textbf{Quantum computation.}
The $\Jcost$-cost bound applies to all variational trajectories, including those that model quantum evolution.
Does the quadratic structure of $\Jcost$ near equilibrium ($\cosh t - 1 \approx t^2/2$) give quantum computation an advantage within the same defect budget?
\end{enumerate}

%% ============================================================
\section{Conclusion}\label{sec:conclusion}
%% ============================================================

We have proved that the $\Jcost$-cost functional imposes absolute bounds on physical computation.
The argument is elementary: non-negative defect, monotone dynamics, and the Archimedean property of the reals are sufficient to show that infinite loops are economically impossible and that every physical computation reaches effective halting within a bound determined by the initial defect and the minimum step cost.

The deepest insight is the unification of logic and computation through cost.
Recognition Science already proved that logical consistency emerges from cost minimization: contradictions are expensive, consistency is cheap, and reality is what is cheap.
This paper extends the same principle to computation: non-halting is expensive, halting is cheap, and reality is what halts.

The universe does not compute by thinking.
It computes by spending defect.
When the defect runs out, the computation is done.

\bigskip

\noindent\textbf{Machine Verification.}\quad
All theorems in this paper have been formalized and verified in Lean~4 using Mathlib, with no \texttt{sorry} placeholders and no custom axioms.
The formalization is in the file \texttt{RecognitionScience/Foundation/AlgorithmicCost.lean}, available at \url{https://github.com/jonwashburn/recognition-science}.

\bigskip

\begin{thebibliography}{99}

\bibitem{RS-LawOfExistence}
J.~Washburn, ``Law of Existence: Rigorous Formalization,'' \texttt{RecognitionScience/Foundation/LawOfExistence.lean}, 2025--2026.

\bibitem{RS-LogicFromCost}
J.~Washburn, ``Logic Emerges from Cost Minimization,'' \texttt{RecognitionScience/Foundation/LogicFromCost.lean}, 2025--2026.

\bibitem{RS-VariationalDynamics}
J.~Washburn, ``Variational Dynamics: The Equation of Motion for the Ledger,'' \texttt{RecognitionScience/Foundation/VariationalDynamics.lean}, 2025--2026.

\bibitem{RS-Thermodynamics}
J.~Washburn, ``Thermodynamics: Temperature, Entropy, and the Canonical Ensemble,'' \texttt{RecognitionScience/Foundation/Thermodynamics.lean}, 2025--2026.

\bibitem{RS-MeasurementMechanism}
J.~Washburn, ``Measurement Mechanism: How Determinism Produces Apparent Randomness,'' \texttt{RecognitionScience/Foundation/MeasurementMechanism.lean}, 2025--2026.

\bibitem{Turing1936}
A.~M.~Turing, ``On computable numbers, with an application to the Entscheidungsproblem,'' \emph{Proceedings of the London Mathematical Society}, 2(42):230--265, 1936.

\bibitem{Levin1973}
L.~A.~Levin, ``Universal sequential search problems,'' \emph{Problems of Information Transmission}, 9(3):265--266, 1973.

\bibitem{Kolmogorov1965}
A.~N.~Kolmogorov, ``Three approaches to the quantitative definition of information,'' \emph{Problems of Information Transmission}, 1(1):1--7, 1965.

\end{thebibliography}

\end{document}
